\section{Positive lower curvature through a perturbed cycle}
\label{art:lower_curvature}

Theorem~\ref{art:local-failure-lower} uses a different regularity balance
from Theorem~\ref{art:local-failure}. Adding a positive quadratic changes
the seed orbit, so exact homogeneity alone no longer supplies a trajectory.
We first prove persistence on a stable manifold and then copy one entire
perturbed orbit at all scales. The first step yields a global failure
from a suitable matrix; the copying step makes bad starts approach the
minimizer of one fixed objective; a final fixed change of coordinates
makes their common initial matrix the identity. The price of retaining
a positive lower curvature bound is the loss of second differentiability
of the branches at the accumulation point. Throughout this appendix the
line-search constants are $c_1=10^{-4}$ and $c_2=1/2$.

\subsection{A strongly convex global example}
A common positive quadratic leaves the active-branch geometry unchanged
and makes every sublevel compact. It perturbs the gradients, however, so
the homogeneous seed is no longer an exact orbit. Close to the corner,
the quadratic gradient is small relative to the fixed slopes. The task
is to turn this smallness into an exact infinite trajectory, with every
rejection and acceptance retained.

Use the fixed slopes of Appendix~\ref{art:seed_certificate} and let
\[
 Q_\mu(x)=\tfrac\mu2\|x\|^2+G(x),\qquad\mu>0.
\]
The vector
\[
 \bar g=\frac{(1,-16)^T}{257}
        =\frac{129}{257}g_0+\frac{128}{257}g_1
\]
is the least-norm point of $\operatorname{conv}\{g_0,g_1,g_2\}$:
$g_0^T\bar g=g_1^T\bar g=\|\bar g\|^2=1/257$, whereas
$g_2^T\bar g>1/257$. Therefore the unique minimizer is
\begin{equation}\label{art:quadratic_minimizer}
 x_*=-\bar g/\mu,\qquad \min Q_\mu=-1/(514\mu).
\end{equation}
At $x_*$ the first two branches tie with distinct gradients, so the
minimizer is nonsmooth. At zero the subdifferential is the nonstationary
triangle of the original seed.

\begin{proposition}[A perturbed bad orbit]\label{art:quadratic_failure}
For every $\mu>0$ and every neighborhood $V$ of zero, some initial state
$(x_0,H_0)$, $H_0\succ0$, produces a deterministic full-BFGS run on
$Q_\mu$ converging to zero, with positive objective values decreasing to
zero. Its accepted lengths repeat $1,1/4,1/16$, and all accepted and
rejected query points are differentiable and belong to $V$.
\end{proposition}

The right object for persistence is the three-step return map after
dividing out the shrinking scale. The homogeneous cycle becomes a fixed
point on the boundary where that scale is zero. We will verify that it
is hyperbolic and that its stable manifold enters the positive-scale
region. This gives exact perturbed orbits; closeness of a finite
floating-point trajectory would not. To prove the proposition, use
normalized coordinates
\[
 x=\tau z,\qquad H=\tau M(b,e),\qquad
 M(b,e)=\begin{pmatrix}1&b\\b&e\end{pmatrix},\qquad\tau>0.
\]
The boundary seed is $P_*=(z_0,0,1,0)$ in coordinates $(z,b,e,\tau)$.
Keep the input $\tau$ fixed during three substeps. Starting with
$(Z,M)=(z,M(b,e))$, for $j=0,1,2$ perform
\begin{align}
 \gamma_j&=g_j+\mu\tau Z,\qquad S=-\alpha_jM\gamma_j,
       \qquad Z^+=Z+S,\label{art:lower_return_substep}\\
 Y&=g_{j+1}-g_j+\mu\tau S,\notag\\
 M^+&=\left(I-\frac{SY^T}{Y^TS}\right)M
          \left(I-\frac{YS^T}{Y^TS}\right)+\frac{SS^T}{Y^TS}.\notag
\end{align}
with cyclic gradient indices and $\alpha=(1,1/4,1/16)$.
For the resulting $(Z_3,M_3)$, let $r=(M_3)_{11}$ and define
\begin{equation}\label{art:lower_return}
 \mathcal T_\mu(z,b,e,\tau)
 =\left(Z_3/r, (M_3)_{12}/r, (M_3)_{22}/r, \tau r\right).
\end{equation}
The seed closure gives $\mathcal T_\mu(P_*)=P_*$ and $r(P_*)=\rho$.
All denominators are nonzero there. Hence the map is real analytic on an
open neighborhood of $P_*$, using negative $\tau$ only to define this
analytic extension.

\begin{lemma}[Implementation and hyperbolicity]\label{art:lower_hyperbolic}
On a sufficiently small neighborhood of $P_*$, positive-$\tau$ states
follow exactly the three query lists of Table~\ref{art:trial_table} under
the actual algorithm, and \eqref{art:lower_return} is their return map.
Its derivative has eigenvalues
\[
 \rho^{-1},\quad\rho^{-1},\quad\lambda_s,\quad\lambda_u,\quad\rho,
 \qquad .038<\lambda_s<.039,\quad1.564<\lambda_u<1.565.
\]
The stable eigenspace contains a vector whose $\tau$ coordinate is one.
\end{lemma}
\begin{proof}
At a query $\tau Z$, the value divided by $\tau$ is
$G(Z)+\mu\tau\|Z\|^2/2$ and the gradient is $g_i+\mu\tau Z$.
The common quadratic does not affect which slope is active. All finite
branch gaps and normalized search comparisons are continuous in the
normalized state. The strict seed margins therefore persist on a common
neighborhood. The physical step and secant are exactly the ones in
\eqref{art:lower_return_substep} after restoring the scale, proving
agreement with the actual update and the entire search tree.

At $\tau=0$ the matrix return is independent of $z$. Let $J$ be the
derivative at $(b,e)=(0,1)$ of its normalized shape map.
The following rational enclosures specify the spectral certificate:
\begin{equation}\label{art:shape_certificate}
 10^{10}J\ \in
 \begin{pmatrix}
 [11123986249,11123986253]&[6459701025,6459701029]\\
 [7515015701,7515015708]&[4905136882,4905136890]
 \end{pmatrix}.
\end{equation}
They are obtained by differentiating the rational substeps
\eqref{art:lower_return_substep} at $\tau=0$, then applying the interval
rules of Appendix~\ref{art:seed_certificate} to
\begin{align*}
 c&\in[318481427627118,318481427627119]/10^{15},\\
 d&\in[-1025232119675749,-1025232119675748]/10^{15}.
\end{align*}
Both intervals contain the defining root box. This is a finite rational
derivative calculation: for a product use
$D(ab)=a\,Db+b\,Da$, and for an inverse use
$D(a^{-1})=-a^{-2}Da$, with denominators bounded away from zero.
Substitution of \eqref{art:shape_certificate} in
$p(\lambda)=\lambda^2-\operatorname{tr}(J)\lambda+\det J$ gives
\[
 p(.038)>0,\qquad p(.039)<0,\qquad
 p(1.564)<0,\qquad p(1.565)>0.
\]
These rational sign comparisons locate two distinct real roots and exhaust
the degree-two polynomial; no numerical eigenvector is used.

At $\tau=0$, $Z_3=z+v(b,e)$ and $r$ is independent of $z$.
The last coordinate of the return is $\tau r$. Thus its full derivative
has the block form
\[
 D\mathcal T_\mu(P_*)=
 \begin{pmatrix}
 \rho^{-1}I_2&C&\mu a\\0&J&\mu b'\\0&0&\rho
 \end{pmatrix}.
\]
The unspecified blocks are finite derivatives of the displayed rational
map and do not affect its spectrum. None of the five eigenvalues is zero
or has modulus one. Moreover $\rho$ is not an eigenvalue of the upper
left four-dimensional block. Assigning last coordinate one and solving
the remaining four linear equations gives a $\rho$-eigenvector. This
proves the last assertion and the lemma.
\end{proof}

\begin{proof}[Proof of Proposition~\ref{art:quadratic_failure}]
The inverse function theorem makes the return map a local diffeomorphism.
The classical local stable-manifold theorem provides a two-dimensional
smooth stable manifold through $P_*$, tangent to the stable eigenspace;
\preprintcite{DY} gives an exposition of this principle and its graph
transform proof. By the last assertion of
Lemma~\ref{art:lower_hyperbolic}, this manifold contains points with
arbitrarily small positive $\tau$.
Choose such a point whose forward return orbit stays in the implementation
neighborhood and tends to $P_*$. Shrinking the neighborhood keeps $r>0$
and $M(b,e)\succ0$, so $\tau^+=\tau r$ preserves positivity. The actual
intermediate matrices remain positive definite by positive curvature.

The physical iterates are $x_{3m}=\tau_mz_m\to0$. Intermediate points
and all queries are $\tau_m$ times bounded continuous functions of the
normalized state, so they tend to zero too. Their normalized objective
values stay positive by the seed's strict positivity. Armijo makes the
values decrease to zero, which exceeds \eqref{art:quadratic_minimizer}.
The initial stable-manifold point can be taken sufficiently close to
$P_*$ that the entire local stable orbit, and hence every physical query,
lies in the prescribed $V$. This proves the proposition. The initial
matrix is supplied by the stable-manifold point, not presumed to be
identity or obtained by rounded forward shooting.
\end{proof}

\subsection{An accumulation embedding with four strongly convex branches}
The previous proposition places a bad orbit near its wrong corner,
whereas the desired local counterexample needs bad starts near the
minimizer. We therefore place shrinking copies around centers $p_n$
converging to that minimizer. The spatial factor will be $a_n$, and the
objective factor $a_n^2$, which preserves the same inverse matrix by
\eqref{art:matched-scaling}. To accommodate the common quadratic, the
three scalar branch slopes are shifted so that the quadratic cross term
restores the original seed gradients inside every copy.

The scalar interpolation must have value $a_n^2$ and derivative $a_n$
at its $n$th node, with an affine neighborhood wide enough to contain
all seed queries after scaling. A monotone cubic ramp between these
plateaux gives a convex function with Lipschitz derivative. The node
slopes divided by the node positions stay of order one, so this
construction need not, and will not, have a second derivative at zero.

Fix
\begin{gather*}
 \mu=1/4096,\qquad c_*=(2,-32)^T,\qquad q_i=g_i-\mu c_*,\\
 b_i=q_i^Tc_*,\qquad r=1/128,\qquad a_n=r^n.
\end{gather*}
Here $b_0=b_1=1791/1024$ and $1.7<b_i<34$. For each of the two distinct
values of $b_i$, set
\[
 t_b=\frac{1-(1+r)/b}{1-r},\qquad\varepsilon=1/1024.
\]
Use the same cubic ramp $P_b$ and primitive $J_b$ as in
\eqref{art:upper_scalar}, with this center and width. For
$ba_{n+1}\le x\le ba_n$ put
\begin{equation}\label{art:lower_scalar}
 \psi_b(x)=(ra_n)^2+ba_n^2(1-r)
       \{r\xi+(1-r)J_b(\xi)\},\qquad
 \xi=\frac{x/(ba_n)-r}{1-r}.
\end{equation}
Set $\psi_b(0)=0$, $\psi_b(x)=x^2/2$ for $x<0$, and
$\psi_b(x)=1+x-b$ for $x\ge b$.

On each positive interval,
\[
 \psi_b'=a_n\{r+(1-r)P_b(\xi)\},\qquad
 \psi_b''=P_b'(\xi)/b,\qquad0\le P_b'\le768.
\]
The identity $r+(1-r)(1-t_b)=(1+r)/b$ matches node values and
derivatives. Their values are $a_n^2$ and $a_n$, respectively.
The positive derivative is monotone and constant near every node.
Its derivative slopes are uniformly bounded, so summing over intervals
and taking a limit at zero proves that $\psi_b$ is globally $C^{1,1}$,
convex and coercive, nonnegative with unique zero at the origin.
The two constant-derivative lengths, divided by $a_n$, are
\[
 b(1-r)(1-t_b-\varepsilon),\qquad
 b(1-r)(t_b-\varepsilon)/r,
\]
both greater than $1/2$ for these values of $b$. Consequently
\begin{equation}\label{art:lower_plateau}
 \psi_b(ba_n+h)=a_n^2+a_nh\qquad(|h|\le a_n/2).
\end{equation}
At zero, the difference quotients of $\psi_b'$ along $ba_n$ and
$(b+1/4)a_n$ equal $1/b$ and $1/(b+1/4)$, respectively. Thus its second
derivative at zero does not exist.

Define
\begin{equation}\label{art:lower_objective}
 \begin{aligned}
 f_i(x)&=\tfrac\mu2\|x\|^2+\psi_{b_i}(q_i^Tx)\qquad(i=0,1,2),\\
 f_3(x)&=\tfrac\mu2\|x\|^2+x_2,\qquad F=\max_i f_i.
 \end{aligned}
\end{equation}
Each branch is globally $\mu$-strongly convex and $C^{1,1}$.
For the first three branches, the gradient-Lipschitz constants are at most
$L_i=\mu+\max\{1,768/b_i\}\|q_i\|^2$. Direct substitution of the
rational root bounds gives $L_0<439$, $L_1<446$ and $L_2<27$; hence
$512$ is a common constant, also covering the fourth branch.
Since $F(x)\ge\mu\|x\|^2/2$ with equality zero at the origin, the
sublevels are compact and the minimizer is unique. The first three
gradients at zero vanish, whereas the fourth is $(0,1)^T$, so the
minimizer is nonsmooth.

At a differentiable point of a finite maximum of $L$-smooth functions,
all active gradients agree. Inactive branches stay inactive locally;
taking the maximum of the active upper Taylor bounds therefore gives
the local upper Taylor bound with constant $L$. Applying this with
$L=512$ and combining it with strong convexity proves
\[
 \mu I\preceq\nabla^2F(x)\preceq512I
 \quad\text{whenever }F\text{ is twice differentiable}.
\]
This does not require a globally Lipschitz gradient for the maximum itself.

\subsection{Copying one entire orbit and fixing the initial matrix}
Let $p_n=a_nc_*$ and $C_*=1+\mu\|c_*\|^2/2$.
For $y$ in a fixed sufficiently small neighborhood of zero,
\eqref{art:lower_plateau} yields
\begin{equation}\label{art:lower_local_model}
 F(p_n+a_ny)=a_n^2\{C_*+Q_\mu(y)\}\qquad(n\ge1).
\end{equation}
Indeed, $|q_i^Ty|\le1/2$ puts each old branch in its affine scalar
region, and expansion of the common quadratic uses $q_i+\mu c_*=g_i$.
The fourth branch is inactive there: its residual $a_n(-32+y_2)$ is
negative, while each old scalar residual is at least $a_n^2/2$.

Choose, by Proposition~\ref{art:quadratic_failure}, one bad orbit
$(\bar x_k,\bar H_k)$ all of whose queries lie in this neighborhood.
For every $n$, start at $p_n+a_n\bar x_0$ with the same matrix
$\bar H_0$. Induction on the actual query tree gives
\[
 x_k^{(n)}=p_n+a_n\bar x_k,\qquad H_k^{(n)}=\bar H_k.
\]
This is \eqref{art:matched-local-model} with spatial factor $a_n$
and objective factor $a_n^2$. The covariance calculation
\eqref{art:matched-scaling} preserves the matrix, while function
differences and directional products acquire the same positive factor.
Thus each rejection, acceptance and updated matrix is preserved. Equation~\eqref{art:lower_local_model} applies to every
query because the whole chosen seed orbit lies in its domain.
Thus $x_k^{(n)}\to p_n$ and
$F(x_k^{(n)})\downarrow a_n^2C_*>0$. The three active gradients at
$p_n$ are $a_ng_i$, whose convex hull excludes zero.
All query norms are at most
$a_n(\|c_*\|+\sup_{\rm seed\ queries}\|\bar x\|)\to0$.

The seed supplied by the stable manifold has a positive-definite initial
matrix $\bar H_0$, common to every copy. Identity initialization is
obtained by changing coordinates once for the whole objective and the
whole sequence of starts, rather than selecting a different objective
for each start. Set $A=\bar H_0^{1/2}$ and $\widetilde F(z)=F(Az)$.
The coordinate identities
\[
 \widetilde g=A^Tg,\qquad \widetilde H=A^{-1}HA^{-T},\qquad
 \widetilde s=A^{-1}s
\]
preserve both scalar line-search tests and
\eqref{art:failure_update}; see also~\cite[Section~2]{LO}.
They give $\widetilde H_0=I$ for this one fixed transformation.
Strong-convexity and smoothness constants become
$\underline\mu=\mu\lambda_{\min}(\bar H_0)$ and
$\overline L=512\lambda_{\max}(\bar H_0)$.
All locality and differentiability statements persist under the
invertible linear map. This proves
Theorem~\ref{art:local-failure-lower}.
